\section{Minimal sampling, optimal recovery amplification, and rigidity}
\label{sec:tpc39-minimality}

The root--of--unity construction uses only half as many measurements as a
full Fourier inversion.  We now show that no phase bank can use fewer.

\subsection{Universal phase measurements}

Let \(z\in\C^\times\).  For a Hilbert column family
\(Z=(Z_{-L},\ldots,Z_L)\), define the phase measurement
\begin{equation}
 \cY_Z(z)=\sum_{k=-L}^{L}z^kZ_k.
 \label{eq:tpc39-general-phase-measurement}
\end{equation}
Given nodes \(z_1,\ldots,z_m\) and coefficients
\(a_1,\ldots,a_m\), the rule
\begin{equation}
 Z_0=\sum_{j=1}^{m}a_j\cY_Z(z_j)
 \label{eq:tpc39-general-recovery-rule}
\end{equation}
holds for every Hilbert space and every column family if and only if
\begin{equation}
 \sum_{j=1}^{m}a_jz_j^k=\one_{k=0}
 \qquad(-L\le k\le L).
 \label{eq:tpc39-moment-system}
\end{equation}

\begin{theorem}[Minimal phase count]
\label{thm:tpc39-minimal-phase-count}
If \eqref{eq:tpc39-general-recovery-rule} holds universally, then
\begin{equation}
 \boxed{m\ge L+1.}
 \label{eq:tpc39-minimal-phase-count}
\end{equation}
This remains true with arbitrary nonzero complex nodes and arbitrary
complex recovery weights.
\end{theorem}

\begin{proof}
Suppose \(m\le L\) and form
\begin{equation}
 Q(z)=\prod_{j=1}^{m}(z-z_j)
 =\sum_{k=0}^{m}c_kz^k.
 \label{eq:tpc39-annihilating-polynomial}
\end{equation}
Its constant coefficient is
\(c_0=(-1)^m\prod_jz_j\ne0\).  Take scalar columns
\(Z_k=c_k\) for \(0\le k\le m\) and zero otherwise.  Then every
measurement in \eqref{eq:tpc39-general-recovery-rule} is \(Q(z_j)=0\),
whereas \(Z_0=c_0\ne0\), a contradiction.
\end{proof}

\begin{corollary}[Optimality of the quotient bank]
\label{cor:tpc39-regular-polygon-minimal}
The nodes
\begin{equation}
 z_\nu=\omega^\nu,
 \qquad0\le\nu<q=L+1,
 \label{eq:tpc39-regular-polygon-nodes}
\end{equation}
with weights \(a_\nu=1/q\) attain the lower bound in
\cref{thm:tpc39-minimal-phase-count}.
\end{corollary}

\begin{proof}
The moment system \eqref{eq:tpc39-moment-system} is additive
orthogonality modulo \(q\), and \(q>|k|\) for every nonzero alias index.
\end{proof}

The lower bound concerns recovery of only \(Z_0\).  Recovering every one
of the \(2L+1\) arbitrary columns requires a measurement matrix of rank
\(2L+1\), and therefore at least \(2L+1\) samples.  The full DFT in
\eqref{eq:tpc39-full-DFT-Parseval} attains that different lower bound.

\subsection{Noise amplification}

Assume now that \(z_j\in\T\).  For a recovery rule satisfying
\eqref{eq:tpc39-moment-system}, define
\begin{equation}
 \cA(a):=m\|a\|_2^2.
 \label{eq:tpc39-noise-amplification}
\end{equation}
The Cauchy bound associated with the rule is
\begin{equation}
 \|Z_0\|^2
 \le\|a\|_2^2\sum_{j=1}^{m}\|\cY_Z(z_j)\|^2.
 \label{eq:tpc39-general-recovery-Cauchy}
\end{equation}

\begin{proposition}[Optimal normalized amplification]
\label{prop:tpc39-optimal-conditioning}
Every universal recovery rule satisfies
\begin{equation}
 \boxed{\cA(a)\ge1.}
 \label{eq:tpc39-amplification-lower-bound}
\end{equation}
The regular polygon rule in
\cref{cor:tpc39-regular-polygon-minimal} has \(\cA=1\).
\end{proposition}

\begin{proof}
The moment at \(k=0\) says \(\sum_ja_j=1\).  Hence
\(1\le m\sum_j|a_j|^2\).  Equal weights \(1/q\) give equality.
\end{proof}

Thus the quotient bank is simultaneously sample--minimal and optimal for
the normalized recovery amplification.  Its estimate
\eqref{eq:tpc39-condition-one-recovery} is the equality case of
\eqref{eq:tpc39-general-recovery-Cauchy} after the natural average
normalization.  The full tomography operator is singular, so this statement
must not be read as a condition number for that operator.

\subsection{Rigidity of positive minimal rules}

The regular polygon is not merely one positive optimal construction.

\begin{theorem}[Positive--rule rigidity]
\label{thm:tpc39-positive-rule-rigidity}
Let \(m=q=L+1\), let \(z_1,\ldots,z_q\in\T\), and let
\(w_1,\ldots,w_q>0\).  Suppose
\begin{equation}
 \sum_{j=1}^{q}w_jz_j^k=\one_{k=0}
 \qquad(-L\le k\le L).
 \label{eq:tpc39-positive-moments}
\end{equation}
Then
\begin{equation}
 w_j=\frac1q\quad(1\le j\le q),
 \label{eq:tpc39-rigid-weights}
\end{equation}
and, after reordering, there is \(\rho\in\T\) such that
\begin{equation}
 \{z_1,\ldots,z_q\}
 =\{\rho\omega^\nu:0\le\nu<q\}.
 \label{eq:tpc39-rigid-nodes}
\end{equation}
\end{theorem}

\begin{proof}
Form the square matrix
\begin{equation}
 U_{j,r}=\sqrt{w_j}\,z_j^r,
 \qquad1\le j\le q,
 \quad0\le r\le q-1.
 \label{eq:tpc39-positive-rule-matrix}
\end{equation}
By \eqref{eq:tpc39-positive-moments},
\((U^*U)_{r,s}=\delta_{r,s}\), since \(|s-r|\le q-1=L\).
Thus \(U\) is unitary.  The squared norm of row \(j\) is \(qw_j\),
which proves \eqref{eq:tpc39-rigid-weights}.  Orthogonality of two
distinct rows gives
\[
 \sum_{r=0}^{q-1}(z_j\overline{z_i})^r=0.
\]
Therefore every quotient \(z_j\overline{z_i}\) is a nontrivial
\(q\)-th root of unity.  The \(q\) distinct nodes are one rotation of all
\(q\)-th roots.
\end{proof}

\subsection{Every minimal twist is necessary}

\begin{corollary}[Deletion no--go]
\label{cor:tpc39-deletion-no-go}
Delete any one measurement from the minimal regular polygon bank.  There
are scalar alias columns with
\begin{equation}
 \cY_Z(z_\nu)=0
 \quad\text{at every retained node},
 \qquad Z_0\ne0.
 \label{eq:tpc39-deletion-witness}
\end{equation}
Hence estimates for all but one minimal twist do not give
coefficient--free exact equality recovery.
\end{corollary}

\begin{proof}
There remain \(L\) nodes.  Use the annihilating polynomial in
\eqref{eq:tpc39-annihilating-polynomial}.
\end{proof}

If an analytic argument has exceptional frequencies, the recovery bank
must therefore be made redundant or the exceptional outputs must be
controlled separately.  Treating an unproved minimal twist as negligible
is not a valid recovery step.
